% s-363-p363-m002; original PDF 363, 364
\Aphorism{s-363-p363-m002}{407}
Wenn wir die Welt, in der wir leben, als Ideal einer causalen, determinirten Welt aufstellen und ihre Ordnung, Verknüpfung, Übereinstimmung, Gesetzmässigkeit vollkommen befriedigend finden, so kann ein Betrachter jenseits von Mensch und Erde achselzuckend sagen: Jedem gefällt seine eigene Nase, Jeder nennt sein Chaos Kosmos. Ist es denn so ausgemacht, dass ein Lebewesen mit höheren Ansprüchen an \Greek{νόμος} und \Greek{ἀνάγκη} sich in unserer Welt heimisch fühlen würde, anstatt ihr vielleicht die Censur „verwirrt und verwirrend bis zur Betäubung“ zu geben? Wir gestehen zu, dass unser \Sperr{Traum} Bilder und Vorstellungen völlig regellos wechselt, meinen dagegen, die Wirklichkeit zeige darin Ruhe und Ordnung. Wohl, sobald wir in der wachen Erinnerung die Traumbilder am ruhigen Hintergrund der Wirklichkeit vorüberziehen lassen, erscheinen sie wechselnd und gesetzlos; im Traume selbst glauben wir so regelrecht zu leben und zu handeln wie nur möglich. Für den Wahnsinnigen ist seine Wahnvorstellung real, seine Ideenflucht ruhig-syllogistische Folge. Hätten wir für die „Wirklichkeit“ noch eine andere Realitätsklasse zum Vergleich, wer weiss, wie verzerrt und regelwidrig sie uns dann erschiene? Ohne diese Controlle von Aussen ist unser immanenter Determinismus eitel Kirchspielweisheit; wir beurtheilen die Erfahrung nach einem Schematismus, den wir (oder unsere biologischen Ahnen) erst aus der Erfahrung gezogen haben, und freuen uns kindlich, wie alles stimmt. Erinnern wir uns doch, wie jung der Determinismus als philosophische \Sperr{Autorität} ist, und vor wie kurzer Zeit man noch Agentien ausserhalb des Causalgesetzes anerkannte. Hier wäre gewiss für jenen Intellect mit höheren Ansprüchen unser Kosmos lange nicht kosmisch genug: eine schöne Ordnung, würde er sagen, die nicht einmal in den sie spiegelnden Köpfen Ordnung schafft! — Diesen Rath, das deterministische Wesen der Welt nicht zu überschätzen, giebt ein Philosoph, der im Übrigen selbst bei höchster Schärfe und Bestimmtheit der \Sperr{inneren} Weltverknüpfung die chaotische Willkürlichkeit des absoluten Seins nicht gefährdet sieht.
\par

